\section*{Chapter \ref{ch:b3:central-banking} --- Central Banking in Practice}

\begin{solution}{exo:b3:central-banking:1}
Federal Reserve: from 6.2~\% to 36.9~\% of GDP, 30.8 points; Eurosystem: from 16.3~\% to
68.1~\%, 51.8 points. By 2025 the Fed's fell by 15.3 points (to 21.6~\%), the
Eurosystem's by 29.2 (to 38.8~\%): the Eurosystem's shrank more in points, and about
as fast in proportion (43~\% against 42~\%). % figdata: cb_assets_points
\end{solution}

\begin{solution}{exo:b3:central-banking:2}
Central bank: assets, government bonds $+100$; liabilities, \omterm{def:b3:central-banking:balance-sheet}{bank reserves} $+100$.
Commercial banks: assets, government bonds $-100$ and reserves $+100$; their
liabilities do not change. Reserves are created by the purchase.
\end{solution}

\begin{solution}{exo:b3:central-banking:3}
Demand gives $r = 3.5 - (R - 20)/40$ between 20 and 60. With 30: $r = 3.25~\%$; with 50: $2.75~\%$;
with 70, beyond satiation: $r = 2.5~\%$, the deposit rate.
\end{solution}

\begin{solution}{exo:b3:central-banking:4}
Futures and overnight-index-swap rates that price rate rises earlier than the
guidance says; long yields that do not fall, or rise, on the day of the guidance;
and surveys of forecasters whose expected path ignores it.
\end{solution}

\begin{solution}{exo:b3:central-banking:5}
$(1 + 2 + 3)/3 + 0.4 = 2.4~\%$. After the guidance, $(1 + 1.5 + 2.5)/3 + 0.4 = 2.07~\%$.
\end{solution}

\begin{solution}{exo:b3:central-banking:6}
On the floor the overnight rate follows the deposit rate up to 2.5~\%, one for one,
with no change in reserves. In a corridor with scarce reserves the rate is set
inside the corridor by reserve supply and demand; raising the deposit rate alone
narrows the corridor and raises the value of holding reserves, so the rate rises,
but by less than half a point unless reserves are near satiation.
\end{solution}

\begin{solution}{exo:b3:central-banking:7}
If the two-year change is all expected short rates and expectations beyond two
years do not move, the expectations part of the ten-year change is $2/10$ of the
two-year change: 2, 2.4 and 0.8 basis points, that is 9~\%, 5~\% and 16~\% of the falls;
the rest is \omterm{def:b3:central-banking:qe}{term premium}. The third day moved little: markets had anticipated
the announcement, and an event study then misses its effect.
\end{solution}

\begin{solution}{exo:b3:central-banking:8}
The tested solution is \code{ex08_inflation_peaks.py}, in the chapter's
\code{exercises/} folder. United States: peak 9.1~\% in June 2022, first below 3~\% in June
2023 (2.97~\%). Euro area: peak 10.6~\% in October 2022, first below 3~\% in October 2023
(2.90~\%). % output: ex08_inflation_peaks
\end{solution}

\begin{solution}{exo:b3:central-banking:9}
Long yields that do not move on purchase announcements beyond what the change in
expected policy rates explains; term premia unaffected by large changes in the
supply of long bonds from other sources (the Treasury's own issuance); and equal
effects on assets bought and assets with no link to them.
\end{solution}

\begin{solution}{exo:b3:central-banking:10}
$r = 3~\%$ for $S \le 20$; $r = 2 + (60 - S)/40$ for $20 \le S \le 60$; $r = 2~\%$ for $S \ge 60$. The centre, 2.5~\%,
needs $S = 40$. On the floor the demand curve is flat: an error in forecasting
demand moves the quantity of reserves banks hold, not the rate.
\end{solution}

\begin{solution}{exo:b3:central-banking:11}
$500 \times 1.5~\% = 7.5$ earned. At 0.5~\%: $7.5 - 2.5 = +5$. At 4~\%: $7.5 - 20 = -12.5$. The loss
reduces what the central bank pays to the treasury, so taxpayers bear it. It
does not stop the bank from operating --- it can run with negative equity --- but
it can become a political constraint on its independence.
\end{solution}

\begin{solution}{exo:b3:central-banking:12}
Buyers with deposits below 20~\% of the price: often first-time, younger and less
wealthy buyers. Their borrowing should fall, and prices in the segments and regions
where they buy should rise more slowly. Compare, before and after the cap,
borrowers or regions with many constrained buyers against those with few, in a
difference-in-differences design.
\end{solution}

\begin{solution}{pb:b3:central-banking:bank}
\textbf{1.} Assets: bonds 60, loans 10, foreign reserves 30, total 100. Liabilities:
currency 50, reserves 40, government deposits 10, total 100. % figdata: bank_years, bank_stats

\textbf{2.} Government deposits $-4$; reserves $+4$ (credited to the contractor's bank).

\textbf{3.} Currency $+2$; reserves $-2$.

\textbf{4.} Loans $+5$ and reserves $+5$, to 45: $r = 3.22 - (45 - 20)/40 = 2.60~\%$, below the
policy rate of 2.72~\%.

\textbf{5.} It controls its securities, its loans (through the facilities' rates) and
its foreign assets; currency is what the public wants to hold and government
deposits are the treasury's, so reserves move with both unless the Bank offsets
them.

\textbf{6.} Only banks hold reserves, at the Bank, to settle payments among
themselves; households and firms hold bank deposits and currency.

\textbf{7.} 2.72~\%: with 40 of reserves the rate is mid-corridor.

\textbf{8.} $3.22 - (30 - 20)/40 = 2.97~\%$.

\textbf{9.} $(60 \times 3.0 + 30 \times 3.0 + 10 \times 2.72 - 40 \times 2.22)/100 = (180 + 90 + 27.2 - 88.8)/100 = 2.08$
billion dalens.

\textbf{10.} Assets: long bonds $+40$; liabilities: reserves $+40$, to 80.

\textbf{11.} Reserves of 80 exceed satiation (60): $r$ falls to the deposit rate, 2.22~\%.

\textbf{12.} With abundant reserves the overnight rate equals the deposit rate;
setting that rate equal to the policy rate sets $r = i$ without adjusting reserves.

\textbf{13.} $3.5 - 2.7 = 0.8$ points.

\textbf{14.} $3.5 - 0.3 = 3.2~\%$.

\textbf{15.} Portfolio balance: with fewer long bonds in private hands, investors who
want them accept a lower premium.

\textbf{16.} $(60 \times 3.0 + 40 \times 2.0 + 30 \times 3.0 + 10i - 80i)/100 = 3.5 - 0.7i$ billion dalens.

\textbf{17.} 2020: $3.5 - 0.7 \times 2.23 = 1.94$; 2021: $3.5 - 0.7 \times 3.11 = 1.32$.

\textbf{18.} Above $i = 5.0~\%$.

\textbf{19.} $1 + 10.18 + 0.5 \times 8.18 = 15.27~\%$.

\textbf{20.} The actual rate was 9.9 points lower. Reasons: inflation driven by
supply shocks expected to pass, policy lags (a rate set for 2022 acts in 2023--24),
and the risk of a recession; the rule is a benchmark, not an instruction.

\textbf{21.} At 5.42~\%: $3.5 - 3.79 = -0.29$ billion. At 15.27~\%: $3.5 - 10.69 = -7.19$ billion, 2.1~\%
of GDP.

\textbf{22.} No: the Bank can operate with losses. A large loss may become a political
argument against its independence, which is why some ask that the treasury cover
it in advance.

\textbf{23.} Futures and swap rates that price high rates through 2024; long-term
inflation expectations from surveys and indexed bonds near 2~\%; long yields that
rise on the announcement day.

\textbf{24.} The rule implies a policy rate of 15.3~\% in 2022, at which the Bank's net
interest income would be $-7.2$ billion dalens (2.1~\% of GDP).
\end{solution}
